% Preamble
\documentclass[11pt, aspectratio=169]{beamer}
	% Packages
	\usepackage[english]{babel}
	\usepackage{lipsum}
	\usepackage{mathptmx}
	\usepackage{xcolor}
	% Themes
	\usefonttheme[onlymath]{serif}
	\usetheme[nofirafonts]{focus}
	\renewcommand{\thempfootnote}{\arabic{mpfootnote}}
	% Cover
	\title{Microeconometrics Using Stata}
	\subtitle{Multinomial models}
	\author{
		Jhon R. Ordoñez
		\footnote{\label{1} National University of San Cristóbal de Huamanga}
		\footnote{\label{2} Faculty of Economic, Administrative and Accounting Sciences}
		\footnote{\label{3} Professional School of Economics}
	}
	\date{\today}
% Body
\begin{document}
	% Cover
	\begin{frame}
		\maketitle
	\end{frame}
	% Outline
	\begin{frame}[t]
		\frametitle{Outline}
		\tableofcontents
	\end{frame}
	% Section 1
	\section{Exercise 1}
		% Slide 1
		\begin{frame}
			\frametitle{Exercise 1}
				\begin{block}{Exercise 1}
					Consider the health-status multinomial example of \textcolor{red!50}{\textbf{section 18.9}}.Refit
					this as a \textbf{MNL} model using the \textcolor{blue}{\textbf{mlogit}} command. Comment on the
					statistical significance of regressors. Obtain the \textbf{MEs} of changes in the
					regressors on the probability of excellent health for the \textbf{MNL} model, and
					compare these to those given in \textcolor{blue}{\textbf{section 18.9.5}} for the ordered logit
					model. Using Bayesian information criterion, which model do you
					prefer for these data - \textbf{MNL} or ordered logit?
				\end{block}
		\end{frame}
	% Section 2
	\section{Exercise 2}
		% Slide 1
		\begin{frame}
			\frametitle{Exercise 2}
			\begin{block}{Exercise 2}
				Consider the CL example of \textcolor{red!50}{\textbf{section 18.5}}. Use \textcolor{blue}{\textbf{mus218hk.dta}}, if
				necessary to create this file as in \textcolor{red!50}{\textbf{section 18.5.1}}. Drop the charter boat
				option as in \textcolor{red!50}{\textbf{section 18.7.5}}, using \textcolor{blue}{\textbf{drop if fishmode=="charter"
				mode==4}}, so we have a three-choice model. Estimate the parameters of
				a \textbf{CL} model with regressors \textcolor{blue}{\textbf{p}} and \textcolor{blue}{\textbf{q}} and income, using the \textcolor{blue}{\textbf{cmclogit}}
				command. What are the \textbf{MEs} on the probability of private boat fishing
				of a $\$10$ increase in the price of private boat fishing, a one-unit change
				in the catch rate from private boat fishing, and a $\$1,000$ increase in
				monthly income? Which model fits these data better-the \textbf{CL} model of
				this question or the \textbf{MNP} model of \textcolor{red!50}{\textbf{section 18.7}}?
			\end{block}
		\end{frame}
	% Section 3
	\section{Exercise 3}
		% Slide 1
		\begin{frame}
			\frametitle{Exercise 3}
			\begin{block}{Exercise 3}
				Continue the previous question, a three-choice model for fishing mode.
				Estimate the parameters of the model by \textbf{NL}, with errors for the utility
				of pier and beach fishing correlated with each other and uncorrelated
				with the error for the utility of private boat fishing. Obtain the \textbf{ME} of a
				change in the price of private boat fishing, adapting the example of
				\textcolor{red!50}{\textbf{section 18.6.6}}.
			\end{block}
		\end{frame}
	% Section 4
	\section{Exercise 4}
		% Slide 1
		\begin{frame}
			\frametitle{Exercise 4}
			\begin{block}{Exercise 4}
				Continue with the fishing data example of \textcolor{red!50}{\textbf{section 18.5}}, with \textcolor{blue}{\textbf{beach}} as
				the base alternative and \textcolor{blue}{\textbf{income}} as the case variable. Instead of fitting
				an alternative-specific conditional logit model, fit the random-parameter logit model of \textcolor{blue}{\textbf{section 18.8}} using the \textcolor{blue}{\textbf{cmmixlogit}} command.
				Is this model a significant improvement on the fixed parameter model?
				Use the \textcolor{blue}{\textbf{margins}} command to fit the \textbf{ME} of income on the choice of
				fishing mode. Compare the role of income in this model with that in
				the model of \textcolor{red!50}{\textbf{section 18.5}}.
			\end{block}
		\end{frame}
	% Section 5
	\section{Exercise 5}
		% Slide 1
		\begin{frame}
			\frametitle{Exercise 5}
			\begin{block}{Exercise 5}
				Consider the health-status multinomial example of \textcolor{red!50}{\textbf{section 18.9}}.
				Estimate the parameters of this model as an ordered probit model using
				the \textcolor{blue}{\textbf{oprobit}} command. Comment on the statistical significance of
				regressors. Obtain the \textbf{MEs} for the predicted probability of excellent
				health for the ordered probit model, and compare these to those given
				in \textcolor{blue}{\textbf{section 18.9.5}} for the ordered logit model. Which model do you
				prefer for these data-ordered probit or ordered logit?
			\end{block}
		\end{frame}
	% References
	\begin{frame}[t,allowframebreaks]
		\frametitle{References}
		\bibliography{biblio.bib}
		% Style
		\bibliographystyle{apalike}
		% Authors
		\nocite{cameron2022-I}
		\nocite{cameron2022-II}
	\end{frame}
	% Cover
	\begin{frame}
		\maketitle
	\end{frame}
\end{document}